\documentclass[12pt]{article}

\usepackage[utf8]{inputenc}
\usepackage[T1]{fontenc}
\usepackage[top=1.25in, bottom=1.25in, left=1.5in, right=1.5in]{geometry}
\usepackage{times}
\usepackage{booktabs}
\usepackage{enumitem}
\usepackage{hyperref}
\usepackage{microtype}
\usepackage{parskip}
\usepackage{titlesec}
\usepackage{xcolor}
\usepackage{mdframed}
\usepackage{array}

\hypersetup{
  colorlinks=true,
  linkcolor=blue!60!black,
  citecolor=blue!60!black,
  urlcolor=blue!60!black
}

% Prominent section headings for quick scanning
\titleformat{\section}{\Large\bfseries}{}{0em}{}[\vspace{2pt}\titlerule\vspace{4pt}]
\titlespacing{\section}{0pt}{20pt}{8pt}
\titleformat{\paragraph}[runin]{\bfseries}{}{0em}{}[.]
\titlespacing{\paragraph}{0pt}{10pt}{6pt}

% Pull-quote box
\newmdenv[
  linewidth=1.5pt,
  linecolor=black!40,
  backgroundcolor=gray!8,
  skipabove=12pt,
  skipbelow=12pt,
  innerleftmargin=14pt,
  innerrightmargin=14pt,
  innertopmargin=10pt,
  innerbottommargin=10pt
]{pullquote}

\title{%
  {\Large\textsc{Policy Brief}}\\[8pt]
  \textbf{Ending Gerrymandering:\\A One-Algorithm Solution}\\[6pt]
  \large Track A, Paper 5
}

\author{Giovanni Della-Libera\\
\small\texttt{giodl@microsoft.com}}

\date{May 2026}

\begin{document}

\maketitle
\thispagestyle{empty}

\begin{pullquote}
\textbf{Bottom line.}
A mathematical algorithm applied to public Census data---with no
political input of any kind---produces congressional districts that are
22\% more compact, closely proportional in partisan representation,
and independently verifiable by any citizen in under ten minutes.
One sentence of federal statute can make gerrymandering structurally
impossible.
\end{pullquote}

\vspace{4pt}

%% 01-problem.tex — The gerrymandering problem

\section{The Problem}

\noindent\textbf{Every ten years, after the Census, the party in power in most
state legislatures redraws congressional district boundaries to ensure
that party wins more seats than the voters it represents would otherwise
produce---a practice called gerrymandering.}
In 2019, the U.S.\ Supreme Court held in \textit{Rucho v.\ Common Cause}
that federal courts are powerless to remedy partisan gerrymandering,
leaving Congress as the only institution that can act~\cite{rucho2019}.
Congress has not acted.
The result: in state after state, mapmakers stretch districts into
grotesque shapes, packing opposition voters into a few districts and
cracking communities across many, to predetermine the outcome of
the next decade of elections before a single vote is cast.

\medskip
\noindent\textit{Rucho} explicitly invited Congress to fix the problem.
This brief proposes how.


%% 02-solution.tex — The algorithmic solution

\section{The Solution}

\noindent\textbf{Replace the politician with an algorithm.}

\medskip
\noindent
The algorithm works like splitting a pizza: make one cut through the
middle of the state, dividing it into two equal-population halves.
Cut each half in half.
Repeat until the state has as many pieces as it has congressional seats.
Every cut is guided by one rule---make it as short as possible---so
districts follow natural geographic features instead of zigzagging to
capture partisan voters.

\medskip
\noindent
The algorithm uses U.S.\ Census population counts and tract boundary
geometries, both of which are public data~\cite{census2020pl}.
It never sees voter registration files, election results, or incumbents'
home addresses.
\textbf{It cannot gerrymander because it lacks the information
required to do so.}
This is not a policy choice that a future legislature could reverse;
it is a structural property of the algorithm.

\medskip
\noindent
The algorithm, called \textit{ApportionRegions}, is a direct extension of
\textit{Huntington-Hill}---the mathematical method Congress already uses to
apportion seats among states, mandated by statute since
1941~\cite{huntingtonHill1941}.
ApportionRegions extends the same principle from apportionment to intrastate
district drawing; the two are related in approach but distinct in statutory
scope.
Extending the principle to boundary design requires one sentence of statute.

\medskip
\noindent
Two independent verifiers running the algorithm on the same Census data
produce byte-identical district assignments, auditable through a
SHA-256 hash chain.
Any citizen can verify the result in 15--20 minutes end-to-end,
including data download.


%% 03-results.tex — Three bullet results

\section{What the Evidence Shows}

Applied to all fifty states across the 2000, 2010, and 2020 Census cycles,
the algorithm produces the following results:

\bigskip

\noindent
\begin{tabular}{@{}>{\bfseries\large}p{0.08\textwidth} p{0.86\textwidth}@{}}
+22\% &
More compact than enacted 2020 maps, measured by the Polsby-Popper
ratio of district area to perimeter squared, averaged across all 435 districts.
Compact districts follow rivers and county lines rather than
twisting through neighborhoods to sort voters.
\textit{(Paper B.2~\cite{dellaLibera2026edgeWeighted})}\\[10pt]

7D / 7R &
North Carolina---the state that has generated more federal gerrymandering
litigation than any other in the past decade---receives a
seven-Democratic, seven-Republican delegation under the algorithm,
closely matching the state's evenly divided electorate.
No partisan data was used.
\textit{(Paper B.11~\cite{dellaLibera2026ncRegions})}\\[10pt]

15--20 min &
Any citizen with a laptop and an internet connection can download
the open-source \texttt{redist} tool, fetch public Census data,
run the Vermont verification fixture, and confirm the result themselves
in 15--20 minutes end-to-end including data download.
Full fifty-state replication takes approximately 18 minutes for one
census year, or two to four hours for all three census years.
\textit{(Replication Package A.4)}\\
\end{tabular}

\bigskip

\noindent Additional findings:
the algorithm produces \textbf{more majority-minority districts}
than enacted plans in states above the 42\% minority-population threshold
identified in Paper D.1~\cite{dellaLibera2026vraCompliance},
and reduces the partisan efficiency gap by \textbf{62\%} compared to
enacted plans~\cite{dellaLibera2026efficiencyGap}---without any
partisan fairness criterion in the algorithm specification.\footnote{%
The algorithm produces near-zero efficiency gaps as a byproduct of geometric
neutrality (mean $|EG|=0.04$ vs.\ $0.08$ for enacted maps).
In some geographically sorted states a small systematic gap persists due to
urban concentration (Paper C.5).}


%% 04-adoption.tex — Three adoption paths

\section{Three Paths to Adoption}

\paragraph{1. Federal statute.}
The cleanest solution requires one sentence added to federal law:
\begin{quote}
\textit{Each State shall apportion its congressional districts using
the ApportionRegions redistricting algorithm (an extension of the
Huntington-Hill apportionment method to intrastate district drawing)
applied to decennial Census redistricting data, subject to requirements
of the Voting Rights Act; provided that deviations from the algorithm's
output necessary to comply with Section~2 of the Voting Rights Act of
1965 are documented in a written Gingles justification.}
\end{quote}
This language is analogous to the statutory mandate for Huntington-Hill
apportionment, which has operated without controversy since 1941.
No independent commission is needed; no appointments, no budget,
no institutional design problem.
The algorithm is the commission.

\paragraph{2. State commission or legislative adoption.}
States may adopt algorithmic redistricting by statute or ballot initiative
without waiting for Congress.
The \texttt{redist} tool already supports any state's seat count
and can be configured to respect existing statutory criteria
(municipal boundaries, county lines, communities of interest).
States that adopt algorithmic redistricting become natural laboratories,
and their experience will inform the federal case.

\paragraph{3. Court-ordered remedy.}
After a court finds a gerrymander under state constitutional standards
(as courts have in Pennsylvania, North Carolina, and New York),
algorithmic plans provide a judicially administrable remedy.
Rather than appointing a special master to draw a map---
an ad hoc, unreviewable process---a court can order the parties
to submit the Census data to \texttt{redist} and adopt the output.
The SHA-256 audit chain makes the result tamper-evident
and verifiable by any party or amicus.


%% 05-contact.tex — Repository, dashboard, contact

\section{Resources and Contact}

\paragraph{Open-source code.}
The full \texttt{redist} source code, including the METIS graph-partitioning
library, is available at:\\
\url{https://github.com/giodl73-repo/REDIST}\\
The code is written in Rust and is freely available for inspection,
audit, and independent replication under an open-source license.

\paragraph{Interactive dashboard.}
District maps for all fifty states and all three census years (2000, 2010, 2020)
are viewable in a browser-based dashboard bundled with the repository
at \texttt{web/dashboard.html}.
No software installation is required.

\paragraph{Research portfolio.}
The full research portfolio---more than forty papers covering
the algorithm, legal compliance, validation, and applications---is
maintained at \texttt{research/} in the same repository.
A non-technical visual guide is in Paper A.3.
The reproducibility package is in Paper A.4.

\paragraph{Contact.}
Questions, comments, and collaboration inquiries are welcome.\\[4pt]
\noindent\textbf{Giovanni Della-Libera}\\
\textit{Principal Software Engineer, Microsoft}\\
\texttt{giodl@microsoft.com}


\vspace{16pt}
\hrule
\vspace{8pt}
{\small
\bibliographystyle{plain}
\bibliography{references}
}

\end{document}
